\section{Related Work} \label{sec:related_work}

Our work is situated at the intersection of weakly supervised learning, label
noise robustness, and graph-based evidence aggregation. We group prior work by
the way it uses the observed candidate labels.

\subsection{Standard and Instance-Dependent PLL}
Partial Label Learning (PLL) centers on label disambiguation: identifying the
unique ground-truth label concealed within a set of candidate labels. Early
work formalized learning from ambiguously labeled images, where each instance is
associated with a candidate label set containing one correct label
\cite{cour2009learning}. Cour et al. introduced the ambiguity degree to analyze
when ambiguous supervision can control true classification error up to an
ambiguity-dependent factor, and proposed a convex surrogate loss for ambiguous
multiclass classification. Later PLL studies explored instance-based
propagation \cite{zhang2015solving}, maximum-margin formulations
\cite{yu2016maximum}, risk-consistent objectives \cite{feng2020provably}, and
weighted losses \cite{wen2021leveraged}. PRODEN progressively identifies true
labels by jointly updating the model and the candidate-supported label
confidence \cite{lv2020progressive}. Deep PLL methods further combine
representation learning with disambiguation, including CRDPLL-style consistency
regularization \cite{wu2022revisiting}. CroSel uses a FIFO memory bank of
historical model predictions and cross selection between two networks to select
confident candidate-supported pseudo-labels \cite{tian2024crosel}. These
methods usually assume that the
true label is contained in the candidate set. Many standard PLL losses keep the
observed candidate membership fixed and use it through the loss, risk
estimator, label weights, or current soft targets rather than rewriting the
candidate source. This early perspective is important for our work: the
observed candidate set is not supervision-free noise, but a weak supervisory
object whose usefulness depends on the ambiguity structure. NSE extends this
view to noisy partial labels by extracting epoch-local active supervision from
the restored native prior, rather than persistently rewriting the candidate
source.

Instance-dependent PLL relaxes the random-noise view of candidate generation.
VALEN associates each instance with a latent label distribution and recovers the
distribution through variational label enhancement \cite{xu2021idpll}. IDGP
further decomposes instance-dependent candidate generation into two sequential
parts: the correct label is generated first, and feature-related incorrect
labels are then selected because of annotation uncertainty \cite{qiao2023idgp}.
POP progressively updates the learning model and purifies candidate label sets,
with the goal of expanding the reliable region \cite{xu2022pop}. These works
show that candidate quality is sample-dependent and that the observed candidate
set can encode structured annotation information. Their main operation is still
candidate-set disambiguation or purification, whereas our method uses this
observation to motivate preserving the native candidate/crowd prior as a
constraint on later evidence fusion.

\subsection{Noisy PLL and Persistent Supervision State}
Noisy Partial Label Learning (NPLL) is more difficult because the true label may
be absent from the observed candidate set. PiCO \cite{wang2022pico} and PiCO+
\cite{wang2022pico+} use prototype-based contrastive learning and robust sample
selection to improve deep PLL under ambiguity and noise. IRNet
\cite{lian2023irnet} detects noisy samples and performs label correction.
UPLLRS \cite{Shi2023Unreliable} recursively separates reliable and unreliable
samples. FREDIS \cite{qiao2023fredis} explicitly combines disambiguation and
refinement by moving incorrect labels out of the candidate set and moving
predicted labels into it. Recent NPLL reconstruction work also separates noisy
samples and reconstructs candidate labels when the true label is missing
\cite{peng2025noise}. ALIM \cite{xu2023alim} reduces the impact of detection
errors by adjusting the importance of the initial candidate set and model
outputs.

These methods provide strong baselines, but they expose a design tension that
cannot be reduced to a single ``source rewrite'' category. Standard PLL usually
assumes the true label is already inside the candidate set; many methods
therefore disambiguate a fixed PL source through losses, risk estimators, label
weights, or current candidate-supported soft targets. NPLL breaks this
assumption because the true label may be absent, so recent methods more often
need correction, refinement, promotion, augmentation, or pseudo-target recovery.
Existing methods differ in where model-induced supervision is stored. Some
methods preserve candidate membership and modify only the loss, risk estimate,
or current label-importance weights; ALIM is a close NPLL example because it
trades off the initial candidate set and model outputs without a hard candidate
rewrite. Some methods maintain soft pseudo-target, label-confidence, or
label-distribution states, such as PRODEN, LWS, PiCO, and VALEN. Others modify
or reconstruct candidate membership, as in POP, FREDIS, IRNet, NPLL
reconstruction, and partial-label augmentation in PALS-style training. UPLLRS
instead carries reliable/unreliable sample-state information. Thus, the relevant
distinction is not whether a method uses model-induced supervision; most strong
methods do. The key question is whether that supervision is carried into later
epochs as a \emph{persistent supervision state}: candidate membership,
reconstructed candidate sets, pseudo-targets, label-importance weights, or
sample-reliability state. When noisy-sample detection or pseudo-label prediction
is wrong, such carry-over can preserve the error through later training. This
motivates an alternative: restore the working prior from the native
weak-supervision record at the beginning of each epoch and use model predictions
only as temporary extraction evidence.

\begin{table*}[t]
\centering
\scriptsize
\setlength{\tabcolsep}{2pt}
\begin{tabular}{p{0.10\textwidth}p{0.15\textwidth}p{0.19\textwidth}p{0.12\textwidth}p{0.22\textwidth}p{0.16\textwidth}}
\toprule
\textbf{Method} & \textbf{Supervision state} & \textbf{Representative update} & \textbf{Carry-over} & \textbf{Potential limitation under carry-over} & \textbf{NSE distinction} \\
\midrule
PRODEN / LWS~\cite{lv2020progressive,wen2021leveraged} & label-confidence / label-weight state & update candidate-supported confidence or weights while keeping membership fixed & soft & mistaken confidence concentration can bias soft targets, but does not hard-edit candidate membership & restores prior and extracts topology-controlled active evidence \\
\hline
CroSel~\cite{tian2024crosel} & history queue / selected pseudo-label set & FIFO memory bank supports cross-network selection of stable, confident candidate labels & selected-set state & selected labels depend on historical prediction stability under the standard PLL assumption that the true label is in the candidate set & queue stability inspires salvage, but NSE uses tri-consensus and restores the native NPLL prior \\
\hline
POP~\cite{xu2022pop} & candidate membership & progressive purification removes low-confidence labels & yes & mistaken deletion may remove useful ambiguous support from later training & no persistent candidate deletion \\
\hline
FREDIS~\cite{qiao2023fredis} & candidate membership & score-gap refinement inserts non-candidates and score-gap disambiguation removes candidates & yes & later disambiguation can remove mistaken refinement labels, but both procedures update the working candidate-set state across fusion rounds & restore native prior before extraction \\
\hline
NPLL reconstruction~\cite{peng2025noise} & reconstructed candidate set & sample separation and KNN/model pseudo-label evidence reconstruct shorter candidate label sets for training & working reconstruction & training optimizes reconstructed membership rather than an extracted active set anchored by restored prior & extract active supervision, not a new candidate state \\
\hline
PALS~\cite{pals2025} & pseudo-label/partial-label augmentation & confident top-1 predictions can augment the next-iteration partial label & partial & model-updated features and predictions can influence subsequent pseudo-labelling and partial-label augmentation & model evidence remains epoch-local and prior-constrained \\
\hline
UPLLRS~\cite{Shi2023Unreliable} & split/training-stage state & recursively separates reliable and unreliable subsets for disambiguation and semi-supervised training & yes & errors in the reliable/unreliable split can affect subsequent disambiguation and semi-supervised training & active eligibility is regenerated each epoch \\
\hline
IRNet~\cite{lian2023irnet} & correction state & candidate/non-candidate score-gap detection followed by non-candidate insertion & yes & correction quality depends on the chosen detection stage and is then carried forward & no correction state is written back \\
\hline
PiCO/PiCO+~\cite{wang2022pico,wang2022pico+} & pseudo-target state & pseudo targets are updated and influence contrastive positive pairs & yes & pseudo-target errors may affect later target construction & pseudo-labels are temporary active targets \\
\hline
ALIM~\cite{xu2023alim} & importance weights / adjusted training targets & recomputes weights from initial candidate set and model outputs & no hard rewrite & close to source preservation, but adjustment happens at prediction-level importance rather than topology-controlled extraction & topology extraction from restored prior \\
\hline
\textbf{NSE} & restored working prior & $\tilde{\Omega}_i^t \leftarrow \Omega_i^0$, then temporary evidence extraction & \textbf{no} & depends on extraction quality and active-set precision & \textbf{epoch-wise source restoration} \\
\bottomrule
\end{tabular}
\caption{Persistent supervision state versus epoch-wise restored extraction in PLL/NPLL. The comparison separates where model-induced supervision is stored: some methods keep candidate membership fixed but update weights or targets, while others carry candidate membership, reconstructed candidate sets, pseudo-targets, label-importance weights, or sample-reliability states across epochs. NSE distinguishes the native source record $\Omega_i^0$ from temporary working evidence and restores the working prior before each extraction stage.}
\label{tab:method_positioning}
\end{table*}

\subsection{Graph-Based Extraction from Weak Supervision}
Graph-based PLL uses the local feature neighborhood as an additional source of
weak supervision. IPAL \cite{zhang2015solving} performs instance-based label
propagation in standard PLL. PALS \cite{pals2025} extends the KNN perspective to
NPLL by exploiting noisy partial labels through weighted nearest-neighbor
pseudo-labelling and label smoothing. These methods indicate that local
geometry can recover useful supervision even when individual candidate sets are
ambiguous.

NSE follows this graph-based direction but changes the role of the KNN graph.
Instead of using a fixed similarity kernel to generate pseudo-labels, NSE uses
Topology-DAES as a supervision extraction kernel. Topology-DAES first evaluates
whether local neighbor votes are consistent with the native weak prior through
masked-entropy reliability, and then adapts the propagation temperature
according to neighborhood entropy. Thus, the KNN graph is not merely a
pseudo-label generator; it becomes a reliability-controlled extractor from an
epoch-wise restored weak prior. Instead of replacing or rewriting the candidate
set, the native candidate/crowd prior constrains the model belief before it
enters the second propagation stage. The reliability ratio $r_i$ performs
sample-wise arbitration between Topology-DAES geometric evidence and temporary
model evidence. PiCO+ further shows that distance to class prototypes can help
distinguish clean and noisy examples under NPLL. NSE uses the same intuition in
a weaker, non-persistent form: class-center scores participate in the
model--KNN--prototype consensus used to salvage active samples, but the
prototype-induced pseudo-target is not written into a later-epoch supervision
state. Reliable selection and salvage affect training eligibility, not the
persistent native candidate/crowd record.
